\section{Punto fijo y comparación con punto flotante}\label{sec:punto-fijo}

\subsection{Formatos}

Se usa la notación Q$m.f$: $m$ bits enteros con signo y $f$ fraccionarios, en complemento a 2. Los modelos de punto fijo trabajan con enteros, redondean mitad hacia arriba en cada corte, $(v + 2^{s-1}) \gg s$, y no saturan: los anchos se derivan de cotas de peor caso y el modelo falla si un valor no entra. La tabla~\ref{tab:formatos} resume los formatos.

\begin{table}[h]
  \centering
  \caption{Formatos de punto fijo por nodo.}\label{tab:formatos}
  \begin{tabular}{llcrl}
    \toprule
    Nodo & Rama & Formato & Bits & Cota \\
    \midrule
    Entrada I/Q & ambas & Q2.0 & 2 & $\{-1, 0, +1\}$ \\
    Coeficientes $h_q$ & ambas & Q1.7 & 8 & $88/128 = 0{,}6875$ \\
    Acumulador & tiempo & Q2.7 & 9 & $\sum|h_q| = 234/128$ \\
    Salida I/Q & ambas & Q2.7 & 9 & $1{,}83$ + error \\
    Banco \texttt{work} & frecuencia & Q7.5 & 12 & 42,5 \\
    Twiddles $W$ & frecuencia & Q2.8 & 10 & $\pm1$ \\
    Espectro $H$ & frecuencia & Q2.8 & 10 & 1,39 \\
    Producto complejo & frecuencia & Q10.13 & 23 & — \\
    \bottomrule
  \end{tabular}
\end{table}

En la rama temporal la entrada vale $-1$, 0 o $+1$, así que cada producto es seleccionar o negar $h_q[k]$: no hace falta multiplicador y el acumulador es exacto. Toda la diferencia con el flotante proviene de cuantizar los coeficientes.

En la rama en frecuencia cada valor intermedio es lineal en la ventana de entrada, $v = \sum_n c_n x_n$, de modo que $|\mathrm{Re}\,v|, |\mathrm{Im}\,v| \le \sum_n (|\mathrm{Re}\,c_n| + |\mathrm{Im}\,c_n|)$. Evaluando esa cota en las nueve etapas (cuatro de FFT, el producto por $H$ y cuatro de IFFT) con las constantes ya cuantizadas, y sumando 1 de margen por redondeo, el máximo es 42,5, en la penúltima etapa de la IFFT. Eso fija 7 bits enteros para \texttt{work}. Un caso de prueba construye la ventana QPSK que lleva cada nodo a su máximo y comprueba que no hay desborde y que la cota no sobra un bit.

\subsection{SQNR}

Se mide contra una única referencia para ambas ramas, el FIR temporal flotante con los coeficientes sin cuantizar, de modo que la SQNR incluye el error de cuantizar $h$, $H$ y $W$:
\begin{equation}
  \mathrm{SQNR} = 10 \log_{10} \frac{\sum_n |y_\mathrm{ref}[n]|^2}{\sum_n |y_\mathrm{fxp}[n] - y_\mathrm{ref}[n]|^2},
\end{equation}
con $y_\mathrm{fxp}$ = salida entera$/2^7$, I y Q juntas como valor complejo y todas las muestras alineadas por índice, sin descarte. Los casos son QPSK 2× con semillas 0 a 4 de 1024 símbolos (2048 muestras cada una); se informa la mínima entre semillas.

\begin{table}[h]
  \centering
  \caption{SQNR mínima (dB) de la rama temporal según los bits fraccionarios de coeficientes (filas) y salida (columnas). Barrido con semillas 0 y 1.}\label{tab:barrido}
  \begin{tabular}{rrrrrr}
    \toprule
    coef \textbackslash{} salida & 5 & 6 & 7 & 8 & 9 \\
    \midrule
    5 & 29,8 & & & & \\
    6 & 32,6 & 42,1 & & & \\
    \textbf{7} & 35,9 & 37,8 & \textbf{44,3} & & \\
    8 & 35,9 & 42,1 & 43,1 & 47,4 & \\
    9 & 36,0 & 43,4 & 50,2 & 50,3 & 57,6 \\
    \bottomrule
  \end{tabular}
\end{table}

\begin{table}[h]
  \centering
  \caption{SQNR mínima (dB) de la rama en frecuencia variando un nodo por vez, con coeficientes y salida en 7 bits fraccionarios. Barrido con semillas 0 y 1.}\label{tab:barrido-frec}
  \begin{tabular}{lrrrrrr}
    \toprule
    Bits fraccionarios & 4 & 5 & 6 & 7 & 8 & 9 \\
    \midrule
    \texttt{work} & 39,6 & \textbf{43,0} & 44,2 & 44,5 & 44,5 & 44,5 \\
    $W$ & 31,0 & 33,6 & 39,2 & 41,4 & \textbf{43,0} & 42,9 \\
    $H$ & 32,2 & 37,4 & 39,5 & 41,0 & \textbf{43,0} & 42,3 \\
    \bottomrule
  \end{tabular}
\end{table}

Con 6 bits fraccionarios en coeficientes y salida también se superan los 40\,dB (42,1\,dB en tiempo), pero se eligieron 7 para dejar unos 3\,dB de margen. Con el formato elegido y las cinco semillas, la SQNR mínima es \textbf{44,2\,dB en tiempo} y \textbf{43,0\,dB en frecuencia}. El error de los coeficientes domina: agregar bits a \texttt{work} no mejora más allá de unos 44,5\,dB.

\paragraph{Correspondencia entre ramas en punto fijo.} Las dos ramas redondean en lugares distintos, así que sus salidas fijas no son idénticas: difieren como máximo en 2\,LSB por componente en todos los casos. Por eso cada RTL se compara palabra por palabra con su propio modelo fijo y no con el de la otra rama.

\subsection{Comparación con distintos conjuntos de datos}

La tabla~\ref{tab:flotante-fijo} compara cada modelo de punto fijo con la referencia flotante en los conjuntos de datos de la sección~\ref{sec:verificacion} y en las cinco semillas QPSK. El error se expresa en LSB de salida ($2^{-7}$) por componente I o Q.

\begin{table}[h]
  \centering
  \caption{Punto fijo contra flotante: error máximo y RMS en LSB y SQNR en dB, con el formato elegido.}\label{tab:flotante-fijo}
  \begin{tabular}{lrrrrrrr}
    \toprule
    & & \multicolumn{3}{c}{Tiempo} & \multicolumn{3}{c}{Frecuencia} \\
    \cmidrule(lr){3-5}\cmidrule(lr){6-8}
    Caso & Muestras & Máx. & RMS & SQNR & Máx. & RMS & SQNR \\
    \midrule
    \texttt{impulso} & 32 & 0,40 & 0,14 & 44,3 & 1,00 & 0,19 & 41,7 \\
    \texttt{ceros} & 32 & 0 & 0 & — & 0 & 0 & — \\
    \texttt{extremos\_tiempo} & 16 & 2,09 & 0,97 & 40,9 & 2,40 & 1,00 & 40,6 \\
    \texttt{extremos\_frecuencia} & 144 & 2,09 & 0,75 & 45,4 & 2,09 & 0,81 & 44,7 \\
    \texttt{bordes} & 202 & 1,05 & 0,48 & 43,9 & 2,05 & 0,61 & 41,8 \\
    QPSK, semillas 0–4 & 5 × 2048 & 1,05 & 0,55 & 44,2–44,4 & 2,05–2,24 & 0,64 & 43,0–43,1 \\
    \bottomrule
  \end{tabular}
\end{table}

\begin{itemize}
  \item \textbf{Ceros:} ambas ramas dan salida exactamente nula; la SQNR no está definida porque la referencia es nula.
  \item \textbf{Impulso:} la salida temporal es $h_q$, así que su error es solo el de cuantizar los coeficientes (máximo 0,40\,LSB). La rama en frecuencia agrega el redondeo de $W$, $H$ y \texttt{work} y llega a 1\,LSB.
  \item \textbf{QPSK:} el resultado es estable entre semillas, con error máximo de 1\,LSB en tiempo y 2\,LSB en frecuencia.
  \item \textbf{Extremos:} es el caso con menor SQNR (40,6\,dB). La entrada vale $\pm1\pm j$ en todas las muestras, sin los ceros del sobremuestreo, y lleva la salida a su pico de 1,84; ahí los errores de cuantización de los ocho coeficientes se suman con el mismo signo. QPSK 2× no alcanza ese caso: su pico es 0,92.
  \item \textbf{Bordes:} con tramas de 1 a 25 muestras la rama en frecuencia queda en 41,8\,dB, por debajo de QPSK largo pero sobre el mínimo.
\end{itemize}

En todos los conjuntos con señal ambas ramas superan los 40\,dB y el error no pasa de 2,5\,LSB; un test del repositorio lo comprueba en cada integración.
